% Fragmento traído con \parte{tablas/tab_T14_4} desde contenido.tex

\begin{tablalafrox}{Título pendiente}{tab_T14_4}%
  {Y Y Y Y Y Y}%
  {\cab{Código} & \cab{Paquete de trabajo} & \cab{Entregable} & \cab{Criterio de aceptación} & \cab{Resp.} & \cab{Período}}
  1.2.1 & Documento de especificación de requerimientos y línea base de alcance de la Etapa 1 (H1) & Especificación de requerimientos aprobada (línea base de alcance). & El 100 \% de los requerimientos de la Etapa 1 tiene módulo y criterio de aceptación, y la línea base está firmada en el mes 2. & ARQ & Meses 1 y 2 (H1). \\
  1.2.2 & Documento de reglas de ruteo capturadas al planificador de rutas & Catálogo de reglas de ruteo validado por don Hugo. & Don Hugo firma que las reglas representan su criterio. Las rutas que el software genera con esas reglas resultan operables cuando se comparan con las suyas. & IMP & Meses 1 a 3. \\
  1.2.3 & Documento de especificación de las interfaces sin documentación del ERP, del WMS de 2013 y de la telemetría del tercero & Especificación de cada interfaz (formato, volumen y horario) y prueba de conexión. & Cada interfaz tiene contrato y una prueba de conexión exitosa en el ambiente de Desarrollo antes del H2. & ARQ & Meses 1 a 4. \\
  1.2.4 & Matriz de trazabilidad: requerimiento, módulo, paquete y prueba (H1) & Matriz de trazabilidad. & Ningún requerimiento ofertado queda sin paquete ni sin prueba. & CAL & Mes 2 (H1). Se actualiza con cada cambio aprobado. \\
  1.2.5 & Documento de línea base de alcance de la Etapa 2: canal moderno y costo de servir (H8) & Línea base de alcance de la Etapa 2. & Está aprobada por la Contraparte Técnica en el mes 14. & ARQ & Meses 13 y 14 (H8). \\
\end{tablalafrox}
